\chapter{Model and Formulation}\label{ch:model}
This chapter derives the proposed model in full for the piano setting. The
additions that the strings-and-winds setting makes to this model are specified
together with its results in Chapter~\ref{ch:phase2}, and the mathematics of
the waveform setting is given in Chapter~\ref{app:phase3}. The chapter
proceeds from the score support and the expressive variables defined on it,
through the score graph and the kernel it induces, to the coupling of the
channels, the side information, the observation model and the posterior.
It closes with the hyperparameter fit and the constrained special cases that
recover the two-stage pipeline. Appendix~\ref{app:notation} summarises the
notation.

\section{Setting}
A score is the known discrete support
\begin{equation}
  \score = \{ s_i \}_{i=1}^{N}, \qquad s_i = (p_i,\, b_i,\, d_i),
  \label{eq:support}
\end{equation}
where $p_i$ is the pitch (MIDI number), $b_i \ge 0$ the onset in score time
(beats), and $d_i > 0$ the notated duration (beats). A performance attaches
to each score event a vector of continuous expressive variables. In the piano
setting this vector is
\begin{equation}
  y_i \;=\; \bigl[\, \tau_i,\; \log r_i,\; v_i \,\bigr] \in \R^{3},
  \label{eq:variables}
\end{equation}
and its entries are defined relative to the monotone score-time warp
$T(\cdot)$, which maps score beats to performance seconds. The warp is read
off the alignment. The local tempo, in seconds per beat, is estimated by
finite differences of performed time against score time along the aligned
beat grid, and $T$ is the cumulative integral of the local tempo over score
time, so that $T(b)$ is the tempo-implied arrival of beat $b$. Because $T$ is
estimated from the same performance's aligned beat grid and then treated as
known, its estimation error is not propagated. This is a limitation, and the
strings-and-winds setting addresses it with a leave-one-out warp
(Chapter~\ref{ch:phase2}). Conditional on $T$, the three channels are
\begin{align}
  \tau_i   &= t_i - T(b_i), \label{eq:tau}\\
  \log r_i &= \log \frac{\delta_i}{\,T(b_i + d_i) - T(b_i)\,}, \label{eq:logr}\\
  v_i      &= \tilde v_i / 127 - \overline{\tilde v/127}. \label{eq:vel}
\end{align}
Here $\tau_i$ is the onset residual in seconds, the performed onset $t_i$
minus the tempo-implied onset $T(b_i)$. The articulation $\log r_i$ is the
log ratio of the performed duration $\delta_i$ to the warp-implied duration,
and the dynamics $v_i$ is the centred, unit-range MIDI velocity
$\tilde v_i$. Conditioning on the warp $T$ is what makes the per-note field
$y$ well defined and jointly Gaussian.

Later settings add channels on the same support. The strings-and-winds
setting adds the intonation $c_i$ in cents, the vibrato parameters and the
fundamental curve $f_0$, and the waveform setting adds the harmonic
amplitudes $a$. Not every variable is available, measurable, or even
meaningful on every instrument, and Table~\ref{tab:variables} records for
each one the instruments on which it is realistic. The piano gives reliable
timing, articulation and dynamics but no performer-controlled intonation,
which is why intonation and vibrato are studied on continuously pitched
instruments. The task is score-informed in the sense that $\score$ is given
and the posterior over $y$ is inferred, in the inverse direction, from
observations generated by the same model that, run forward, synthesises a
performance.

{\footnotesize
\begin{longtable}{@{}p{3.1cm}p{2.0cm}p{2.0cm}p{0.9cm}p{4.4cm}@{}}
\caption{Performance variables by setting (P = piano, SW = strings and
winds), with what they are realistic for. Waveform-setting quantities,
the harmonic amplitudes $\mathbf a_i$ and the amplitude envelope $A_i(t)$, are
defined in Chapter~\ref{app:phase3}.}
\label{tab:variables}\\
\toprule
Variable & Symbol & Unit & Setting & Realistic for \\
\midrule
\endfirsthead
\toprule
Variable & Symbol & Unit & Setting & Realistic for \\
\midrule
\endhead
\bottomrule
\endfoot
Score-time warp (local tempo) & $T(\cdot)$ & s/beat, integrated & P & all (defines the score$\to$time warp) \\
Note onset residual & $\tau_i$ & s (report ms) & P & all \\
Articulation / duration ratio & $\log r_i$ & dimensionless & P & all (esp.\ piano) \\
Velocity / dynamics & $v_i$ & norm.\ MIDI & P & keyboard reliable; others approximate \\
Intonation deviation & $c_i$ & cents & SW & voice, strings, winds --- \emph{not} piano \\
Instantaneous $f_0$ curve & $f_{0,i}(t)$ & Hz & SW & monophonic pitched \\
Vibrato (rate, extent, delay) & $f_i^{\mathrm{vib}}, \gamma_i, d_i^{\mathrm{vib}}$ & Hz, cents, s & SW & voice, strings, winds \\
Loudness (audio) & $\ell_i$ & log RMS & SW & monophonic pitched \\
Posterior uncertainty & $\sigma_i$ / CI & same as variable & all & all \\
\end{longtable}}

\section{The score graph}\label{sec:scoregraph}
Let $\graph = (V, E)$ be a weighted graph on the $N$ score events with
Gaussian-kernel edge weights
\begin{equation}
  W_{ij} \;=\; \exp\!\Bigl( -\tfrac{(b_i - b_j)^2}{2\,\ell_b^2} - \tfrac{(p_i - p_j)^2}{2\,\ell_p^2} \Bigr)
  \;+\; w_v\, \mathbf{1}[\text{voice}_i = \text{voice}_j],
  \qquad W_{ii} = 0,
  \label{eq:adjacency}
\end{equation}
with length scales $\ell_b$ (beats) and $\ell_p$ (semitones) and an optional
same-voice bonus $w_v \ge 0$. Throughout, these are untuned defaults,
$\ell_b = 2$, $\ell_p = 4$ and $w_v = 0$. The ASAP corpus carries no voice
labels, so the voice bonus is inactive in every reported experiment and is
retained for corpora that annotate voices. The combinatorial graph Laplacian
is
\begin{equation}
  \Lg \;=\; D - W, \qquad D = \diag\Bigl(\textstyle\sum_j W_{ij}\Bigr).
  \label{eq:laplacian}
\end{equation}
Two baselines are defined on the same support. The temporal-only baseline
replaces $W$ by the chain adjacency of score-time order, so that each note is
coupled only to its temporal neighbours, which gives a smoother of AR(1)
type. The independent baseline drops all edges, so that
$\Lg = 0$ and the notes are a priori unrelated.

\subsection{Choice of Laplacian}
Alongside the combinatorial $\Lg = D - W$ we consider the symmetric
normalized Laplacian
\begin{equation}
  \Lg^{\mathrm{sym}} \;=\; I - D^{-1/2} W D^{-1/2},
  \label{eq:lapsym}
\end{equation}
which rescales each note by its degree. The choice matters in principle
because the score graph is markedly irregular. A note inside a dense chord
has many more neighbours than one in a sparse melodic line, so without
normalisation the prior couples chordal notes far more tightly than melodic
ones. Which Laplacian is used is part of the first of the two design axes
described at the end of this section.

\subsection{Music-theoretic graph constructions}\label{sec:graphvariants}
The base adjacency \eqref{eq:adjacency} treats pitch as a ruler, since
proximity is measured in semitones. Music theory suggests two different, and
logically distinct, interventions. One \emph{replaces} that ruler, the other
\emph{adds} relations on top of it. Both are evaluated in
\S\ref{sec:kernels}, and the distinction between them is what makes that
result interpretable.

The first intervention replaces the semitone ruler by a tonal pitch metric.
Write $\mathrm{pc}(p) = p \bmod 12$ for the pitch class and place it on the
circle of fifths via $\pi(p) = 7\,\mathrm{pc}(p) \bmod 12$. The
circle-of-fifths distance is the shorter way around that ring,
\begin{equation}
  d_5(i,j) \;=\; \min\bigl(\,|\pi(p_i) - \pi(p_j)|,\; 12 - |\pi(p_i) - \pi(p_j)|\,\bigr)
  \;\in\; \{0,\dots,6\},
  \label{eq:fifths}
\end{equation}
so a perfect fifth is $1$ step, a major second $2$, and a semitone $5$. This
is tonal proximity, and it is by design unlike semitone proximity.
Substituting it for the pitch term of \eqref{eq:adjacency}, dropping the
voice bonus and adding an octave-band locality term alongside, gives
\begin{equation}
  W^{\mathrm{tonal}}_{ij} \;=\; \exp\!\Bigl( -\tfrac{(b_i - b_j)^2}{2\ell_b^2}
    \;-\; \tfrac{d_5(i,j)^2}{2\ell_5^2}
    \;-\; \tfrac{\bigl((p_i - p_j)/12\bigr)^2}{2\ell_{\mathrm{oct}}^2} \Bigr),
  \qquad W^{\mathrm{tonal}}_{ii} = 0 .
  \label{eq:tonaladj}
\end{equation}
The third term is necessary rather than decorative. Since $d_5$ is blind to
register, notes four octaves apart would be maximally close without an
octave-displacement penalty, whose length scale $\ell_{\mathrm{oct}}$ is
measured in octaves. The score-time term is unchanged.

The second intervention keeps the base graph and adds two families of edges.
A third family, which links notes whose sustained durations overlap, was
evaluated on the development set and is not used. The augmented adjacency is
\begin{equation}
  W^{\mathrm{chord}}_{ij} \;=\; W_{ij}
    \;+\; w_{\mathrm{ch}}\, \mathbf{1}\bigl[\, b_i = b_j \,\bigr]
    \;+\; w_{\mathrm{vl}}\, \mathbf{1}\bigl[\, b_i \neq b_j,\;
        |b_i - b_j| \le \Delta_{\mathrm{vl}},\;
        1 \le |p_i - p_j| \le 2 \,\bigr],
  \label{eq:harmadj}
\end{equation}
with $\Delta_{\mathrm{vl}} = 2$ beats. The first family links notes struck
together. The second links notes a melodic step apart, one or two semitones,
within a short window, which is stepwise voice leading.

The chord edges should be understood for what they are. The indicator
$\mathbf{1}[b_i = b_j]$ is onset simultaneity, used as a proxy for same-chord
membership. No harmonic analysis is performed, no chord is labelled, and no
key is inferred. Onset simultaneity is the strongest harmonic-simultaneity
signal obtainable from the score alone, and it is also the weakest possible
music-theoretic commitment. Cadential and functional edges, which would
encode tonic and dominant relations, do require a harmonic analysis and are
not implemented (\S\ref{sec:musictheory}). Each family can be removed
independently by setting its weight to zero. The edge weights
$w_{\mathrm{ch}}, w_{\mathrm{vl}}$ are untuned defaults, and a sensitivity
analysis over six settings, scaling each family's weight by $\times 3$ and
by $\div 3$ (\S\ref{sec:kernels}), shows that the resulting gain does not
depend on a narrow range of either weight.

The graph augmented by \eqref{eq:harmadj} is called the \textbf{chord graph},
and its two families the chord and voice-leading edges. They are edges of
harmony in the musical sense, recording which notes sound together and which
move by step. The name avoids the word ``harmonic'', which this thesis
reserves for the harmonic amplitudes $a$ of the waveform synthesiser
(Chapter~\ref{app:phase3}), the overtones of a single note's fundamental. The
two senses never interact.

Together with \S\ref{sec:graphgp}, the choices above fix the graph kernel
completely, and the kernel has two independent design axes. The first is the
graph $\Lg$ itself, which edges exist, under which pitch metric, and which
Laplacian, and it is fixed in this section. The second is the spectral
profile $g$, how sharply high graph frequencies are suppressed, and it is
fixed in \S\ref{sec:graphgp}. The kernel comparison of \S\ref{sec:kernels}
varies both and finds that they matter to very different degrees.

\section{The graph kernel}\label{sec:graphgp}
Each expressive channel is a scalar field over the $N$ notes, and the graph
turns it into a Gaussian process on the score graph, a graph GP, that is, a
multivariate Gaussian whose covariance is built from the graph. Let
$\Lg = U \diag(\nu)\, U^\top$ be the Laplacian's eigendecomposition, with
eigenvalues $0 = \nu_1 \le \nu_2 \le \dots$ playing the role of graph
frequencies. A graph kernel is obtained by shaping this spectrum,
\begin{equation}
  K_G(s) \;=\; U \operatorname{diag}\bigl(g(\nu;\, s)\bigr)\, U^\top,
  \label{eq:spectral}
\end{equation}
with a decreasing spectral profile $g$. Components that vary smoothly along
the graph, those with small $\nu$, keep their prior weight, rapidly varying
components are suppressed, and a single shape parameter $s$ sets how
aggressively. The profile is normalised so that $g(0) = 1$, and $K_G$
therefore carries shape only, with all scale living in the coregionalization
matrix of \S\ref{sec:coreg}. This normalisation is forced rather than
conventional. In the product $B \otimes K_G$ the split of scale between the
two factors is unidentifiable, because $(cB) \otimes (K_G/c)$ is the same
covariance, so the scale must live in exactly one place. Three classical
profiles are used:
\begin{align}
  \text{additive / regularized Laplacian:}\quad & g(\nu;\, s) = \frac{1}{1 + s\,\nu},
  \label{eq:additive}\\[2pt]
  \text{graph Mat\'ern:}\quad & g(\nu;\, s) = \Bigl(\frac{s}{\,s + \nu\,}\Bigr)^{\!\alpha},
  \qquad \alpha \in \mathbb{Z}_{\ge 1},
  \label{eq:matern}\\[2pt]
  \text{diffusion / heat:}\quad & g(\nu;\, s) = e^{-s\,\nu}.
  \label{eq:diffusion}
\end{align}
The Mat\'ern profile is the graph Mat\'ern Gaussian process of Borovitskiy
et al.\ \cite{borovitskiy2021}, with covariance eigenvalues
$\sigma_g^{2}(\kappa^2 + \nu_j)^{-\alpha}$ on the Laplacian eigenvalue
$\nu_j$. Here $s = \kappa^2$, and the scale
$g(0) = \sigma_g^2 \kappa^{-2\alpha}$ is carried by $B$. It is the graph
analogue of a Euclidean Mat\'ern GP through the SPDE link of Lindgren et
al.\ \cite{lindgren2011}. The diffusion profile is its $\alpha \to \infty$
limit under the joint reparameterization $s = \alpha/t$, whereas at fixed $s$
the limit degenerates to zero. The additive profile is not a fourth family.
It \emph{is} the graph Mat\'ern at $\alpha = 1$ with the shape parameter
inverted, since $1/(1+s\nu) = s'/(s'+\nu)$ at $s' = 1/s$, and it is also
known as the regularized-Laplacian kernel \cite{smolakondor2003}.

Under any of these profiles, neighbouring notes are encouraged, not forced,
to share expressive behaviour, and distant notes are coupled only weakly.
The constant profile $g \equiv 1$ gives $K_G = I$, the graph switched off
and the notes a priori independent, which is the ablation that isolates the
graph's marginal value. The additive profile is the default throughout, and
the kernel comparison of \S\ref{sec:kernels} shows that the classical
alternatives perform no differently. Throughout, $\alpha$ is only the
Mat\'ern exponent and $\sigma_g$ the Mat\'ern prior scale, while plain
$\sigma$ is reserved for the posterior standard deviation
(\S\ref{sec:gpfirst}).

\section{Coupling the channels}\label{sec:coreg}
Timing, articulation, and velocity are not independent, since a stressed
note tends to be louder, longer, and delayed together. The model therefore
couples the channels with an intrinsic coregionalization structure
\cite{bonilla2008}, one latent field over the $3N$ (note, channel) pairs with
\begin{equation}
  \operatorname{Cov}\bigl[\, y_{i,c},\; y_{j,c'} \bigr]
  \;=\; B_{cc'}\; \bigl[K_G(s)\bigr]_{ij},
  \qquad \text{i.e.} \qquad
  \operatorname{Cov}[\, y \,] \;=\; B \otimes K_G(s),
  \label{eq:icm}
\end{equation}
where $B \in \R^{3 \times 3}$ is a positive semidefinite matrix learned
through its log-Cholesky factor. Because the $g(0) = 1$ normalisation
resolves the scale ambiguity of \S\ref{sec:graphgp}, $B$ carries all of the
scale. The per-note prior variance of channel $c$ is $B_{cc}\,[K_G]_{ii}$,
since the normalisation fixes the shape of the spectrum rather than the
diagonal of $K_G$, and the off-diagonal entries of $B$ let an observed
channel inform an unobserved one at a coupled note. Constraining $B$ to be
diagonal severs exactly this coupling and makes the three channels
independent scalar graph GPs, which is one of the constraint switches of
\S\ref{sec:special-cases}.

\section{Side information as kernels}\label{sec:features}
Two per-note feature sets summarise what is known about each note beyond the
graph. The first is a set of hand-built score features $X_{\mathrm{feat}}$,
$\sim$25 interpretable score-only columns plus a bias, covering the metric
position in the bar, pitch and melodic interval, contour, local rhythm,
chord and harmonic context, note density, and related quantities. The second
is the set of music-model embeddings $X_{\mathrm{emb}}$, the per-note
embeddings $h_i \in \R^{d}$ of the music model. Both sets enter the model in
the same way, and they differ only in what they encode.

The embedding of a note is read out before the model has seen that note's
own target. The embedding $h_i$ is read at the token preceding note $i$'s
velocity token, which is its duration token. A causal model's hidden state
at position $t$ already contains token $t$, so reading at the velocity token
would hand each note its own velocity, which is the target. The general rule
is that a read-out used at observed notes must never contain its own target.
A bidirectional model instead needs a leave-one-out read-out, in which each
note is read at a masked copy of its own velocity position. Under the strict
protocol, the velocities of held-out notes are also removed from the music
model's input, so the embedding of a held-out note depends only on the score
and on the observed notes' performances.

Each feature set enters the GP as a linear kernel with per-channel scales
$c_f \in \R^{3}_{\ge 0}$, contributing $\diag(c_f) \otimes X_f X_f^\top$ to
the covariance. Each feature matrix has its columns z-scored within the
piece, with a bias column appended, before entering the kernel. Because the
single per-channel scale $c_{f,c}$ is isotropic over columns, this
standardisation affects the model and is not a cosmetic step. The
construction is not a heuristic. A GP whose mean is a linear model
$X_f \beta_c$ with unknown per-channel weights
$\beta_c \sim \N(0,\, c_{f,c} I)$ becomes, once each $\beta_c$ is integrated
out, a zero-mean GP with $c_{f,c}\, X_f X_f^\top$ added to channel $c$'s
kernel. A linear kernel \emph{is} the marginalized Bayesian linear mean. The
feature term is channel-diagonal, so the coregionalization $B$ acts on the
graph component only and cross-channel information flows exclusively through
$B \otimes K_G$. The regression weights are thereby re-inferred for every
piece, jointly with everything else, and their uncertainty flows into the
predictive variance.

The marginalized construction differs from the classical plug-in
alternative, which fixes the weights once. That alternative fits a point
estimate $\hat\beta$ across pieces, subtracts the resulting plug-in mean
$\mu = X_f \hat\beta$, and models the residual with the graph GP. It is one
of the constraint switches of \S\ref{sec:special-cases}, and the attribution
experiments locate most of the proposed model's accuracy gain precisely at
this switch (\S\ref{sec:attribution}).

\section{Observation model and posterior}
\label{sec:gpfirst}
Assembling \S\ref{sec:graphgp} to \S\ref{sec:features}, the latent field is
observed through independent Gaussian noise,
\begin{equation}
  \tilde y \;=\; y + e, \qquad e \sim \N\bigl(0,\, \Sigma_e\bigr),
  \qquad \Sigma_e = \diag(\varsigma^2) \otimes I,
  \label{eq:obs}
\end{equation}
with one learned noise variance $\varsigma_c^2$ per channel. The covariance
of the stacked observation vector, which states the whole model in one
equation, is
\begin{equation}
K \;=\;
\underbrace{B \otimes K_G(s)}_{\text{coupled graph GP}}
\;+\; \underbrace{\sum_f \operatorname{diag}(c_f) \otimes X_f X_f^\top}_{\text{marginalized linear means}}
\;+\; \underbrace{\operatorname{diag}(\varsigma^2) \otimes I}_{\text{observation noise}} .
\label{eq:gpfirst}
\end{equation}
Masked (held-out) notes contribute no likelihood term. Write $o$ for the
observed entries. A held-out note hides all three of its channels, which
keeps the observed block Kronecker-compatible, and exact conjugate GP
regression then gives the posterior over the whole field in closed form:
\begin{equation}
  p(y \mid \tilde y_o) = \N\bigl(m,\, \Sy\bigr), \qquad
  m \;=\; C_{\cdot o}\, K_{oo}^{-1}\, \tilde y_o, \qquad
  \Sy \;=\; C \,-\, C_{\cdot o}\, K_{oo}^{-1}\, C_{o \cdot},
  \label{eq:posterior}
\end{equation}
where $C$ is the latent covariance, the first two terms of
\eqref{eq:gpfirst}, and $K_{oo}$ the observed block of $K$, noise included.
Masked notes are predicted from the observed ones through the conditional
Gaussian, which is the held-out imputation task. The per-note posterior
standard deviation is $\sigma_i = \sqrt{(\Sy)_{ii}}$.

By construction the latent field is a sum of independent Gaussian
components, the coupled graph GP and one draw per feature kernel, so this
posterior decomposes exactly by prior component. The derivation is given at
its single point of use, the head of the attribution analysis
(\S\ref{sec:attribution}). Figure~\ref{fig:arch-phase1} lays the assembly
out as a computation graph. The graph-structure branch and the evidence
branch are built independently from the score, fuse as a sum of Kronecker
terms into the joint kernel \eqref{eq:gpfirst}, and meet the data only in
the likelihood, which is the one block that changes in the later settings
(Chapter~\ref{ch:future}).

\figph{arch_phase1.pdf}{1.0}{The proposed model as a computation graph in the
piano setting. Blue: the graph-structure branch, from the score graph
through $L_{\mathcal G}$ and the spectral shaping $g(\nu;s)$ to $K_G$.
Amber: the evidence branch, in which score features and music-model
embeddings enter as linear kernels, the marginalized Bayesian linear means
of \S\ref{sec:features}. Green: their Kronecker lifts sum with the noise
term into the joint kernel $K$ of Eq.~\eqref{eq:gpfirst}. The likelihood
block (vermillion) is the only element that differs between settings
(cf.\ Figures~\ref{fig:arch-phase2} and~\ref{fig:arch-phase3}). The exact
conjugate posterior returns per-note point estimates with calibrated error
bars.\label{fig:arch-phase1}}

\subsection{Predictive variance for held-out observations}
A held-out observation at entry $i$ is $\tilde y_i = y_i + e_i$, so its
predictive variance must include the observation-noise variance of that
entry. This variance is written $\sigma_e^2$ throughout. In the present
model it is the channel's $\varsigma_c^2$, and in the per-channel form of
Appendix~\ref{app:lineage} it is the per-note $(\Sigma_e)_{ii}$. The
predictive variance is therefore
\begin{equation}
  \operatorname{Var}\bigl[\tilde y_i \mid \tilde y_{o}\bigr]
  \;=\; (\Sy)_{ii} + \sigma_e^2 .
  \label{eq:floor}
\end{equation}
The latent variance $(\Sy)_{ii}$ collapses toward zero at well-determined
notes, so scoring held-out data against $(\Sy)_{ii}$ alone is severely
overconfident, and both NLL and coverage deteriorate sharply.
Equation~\eqref{eq:floor} is therefore required for all held-out
evaluation, and it is referred to below as the variance floor.

\section{Hyperparameters}
\label{sec:eb}
Every hyperparameter is learned jointly, per piece, by maximising the exact
log marginal likelihood, the evidence, of the observed block. The
hyperparameters are the graph shape $s$, the coregionalization $B$, the
feature scales $c_f$, and the per-channel noise $\varsigma^2$, and the
objective is
\begin{equation}
  \log p(\tilde y_o \mid \theta)
  \;=\; \log \N\bigl(\tilde y_o;\; 0,\; K_{oo}(\theta)\bigr),
  \qquad \theta = \bigl(s,\, B,\, \{c_f\},\, \varsigma^2\bigr),
  \label{eq:marglik}
\end{equation}
The objective is optimised in log space for $s$, $c_f$ and $\varsigma^2$ and
in log-Cholesky space for $B$, so positivity and positive semidefiniteness
hold by construction. The optimiser is L-BFGS-B on those unconstrained
coordinates, followed by a derivative-free Nelder--Mead
refinement.\footnote{Where the gradient-based optimiser of SciPy is
unavailable, the fit uses Nelder--Mead alone.} A non-finite objective value
is returned as a large finite value rather than raised as an error, so an
unsuccessful Cholesky factorisation cannot disrupt the line search. One
evidence per piece decides every parameter, including how much to trust each
information source, a choice that a plug-in construction must make once,
across pieces, by cross-validation. A calibration-split alternative
objective instead splits the observed notes and picks $\theta$ minimising
the held-out predictive NLL of the calibration subset under the floored
predictive variance \eqref{eq:floor}.

\subsection{Noise floor}
On a minority of mask realisations the unconstrained maximiser of
\eqref{eq:marglik} degenerates, with $\sigma_e^2 \to 0$, and produces an
overconfident posterior. The same principle as \eqref{eq:floor} therefore
applies to the fit itself. Each noise variance is constrained above a small
floor $\varepsilon_{\mathrm{nv}}$ inside the optimisation, for example a
fixed fraction of the observed residual variance. The constraint binds only
where the unconstrained maximiser would fall below the floor.

\subsection{Validation of the fitted hyperparameters}
Rarely, the maximiser of \eqref{eq:marglik} lands on a degenerate
configuration that the noise floor cannot catch. The fitted field then
extrapolates badly, or reports far too tight intervals, on notes it did not
fit. A safeguard therefore screens the fit on data it did not use. A
calibration subset of the observed notes ($\sim$30\%) is held back, the
fitted GP is conditioned on the remainder, and the fit is rejected if, on
the calibration subset, its RMSE exceeds a small multiple ($2\times$) of the
mean-only RMSE or its NLL exceeds the mean-only NLL by more than $0.5$ nats,
the latter being the overconfidence check. On rejection, the fit falls back
through a fixed sequence: first a refit with the noise floor raised
$5\times$, then a conservative decoupled fit with the graph and feature
kernels off and $B$ diagonal at the observed per-channel variance, which
predicts roughly the mean with a conservative scale. On healthy fits the
screen never fires, so runs with and without the safeguard are identical
except at the degenerate piece-seed combinations.\footnote{The behaviour of
this screen on the piano evaluation, including the one failure mode it
cannot detect, is reported in \S\ref{sec:guard-gp}.}

\section{Constraints and special cases}\label{sec:special-cases}
The model \eqref{eq:gpfirst} has a small set of independent
\textbf{constraint switches}, each of which turns one modelling decision
off:
\begin{enumerate}[leftmargin=2em,noitemsep]
  \item decouple the channels, by constraining $B$ to be diagonal
        (\S\ref{sec:coreg}),
  \item drop the side information, by setting the feature scales $c_f = 0$,
        which leaves a zero-mean graph smoother,
  \item plug the mean in, by replacing the marginalized per-piece weights
        with a cross-piece point estimate $\hat\beta$ and modelling the
        residual (\S\ref{sec:features}),
  \item fix the graph, by holding the graph parameters at constants per
        channel instead of learning them inside the joint evidence.
\end{enumerate}
Turning all four switches on simultaneously yields the classical
\textbf{two-stage pipeline}, a point-estimate trend plus per-channel
residual kriging on the graph, which is the standard construction in the
spatial-statistics and graph-GP literature. This corner of the constraint
space is retained for a methodological reason. It is the null model against
which every hypothesis of the form ``relaxing constraint X helps'' is
tested, so that a gain over it is a statement about the standard
construction rather than about an earlier version of this work. In one
regime, rendering from a short excerpt, where per-piece adaptation has too
little to adapt to, it also remains the better tool (\S\ref{sec:boundary}).
Its per-channel precision-side form, with $\Qg = \lambda I + \eta \Lg$ per
channel as the sparse-GMRF computational route, and its plug-in-mean choices
are recorded in Appendix~\ref{app:lineage}.\footnote{The constrained
systems' posterior formulas and the development lineage tables are kept in
the repository's design record, \texttt{docs/graphgp\_first\_design.md}.}

\subsection{Exact nesting of the two-stage pipeline}
Constraining $B$ to be diagonal, dropping the feature kernels, and fixing
the graph recovers the two-stage pipeline's per-channel model
(Appendix~\ref{app:lineage}) exactly on the shared-shape slice. The joint
marginal likelihood \eqref{eq:marglik} equals the sum of the per-channel
ones under $(\lambda_c, \eta_c) = (1/B_{cc},\, s/B_{cc})$, and the equality
is verified numerically. The per-channel pipeline, fit independently, may
select three different shapes, so the nesting statement is about the model
family and not about coinciding optima. The constrained systems are
therefore not merely analogous baselines but points in the proposed model's
own parameter space, which is what makes the development lineage an
ablation of one model rather than a comparison of two. Re-inserting the
two-stage pipeline's own cross-piece mean as a fixed mean inside
\eqref{eq:gpfirst} reproduces its results almost unchanged. This is how the
proposed model's gain is attributed to the per-piece Bayesian feature
weighting rather than to coupling or to joint fitting as such
(\S\ref{sec:attribution}).

The proposed model and its constrained special cases are evaluated under
one protocol. Predictions on held-out notes are scored on both axes,
recovery and calibration, by the metrics of \S\ref{sec:metrics}, with every
system contrast paired at the level of the piece. The protocol itself,
including the division of the corpus into a development set and a held-out
test set, is described in \S\ref{sec:hygiene}.

% =====================================================================
